\section{Error Troubleshooting}
\label{Error Troubleshooting}

This section records a code-based investigation of intermittent failures during
automated sequences. It focuses on reports that the system appeared to shut
off without an obvious error after several runs, sometimes around 2--4 AM.
The analysis uses the current Python source and the method logs included in
the project. It identifies plausible software paths, not a confirmed
electrical or hardware root cause.

\subsection{Summary of the leading explanation}

The strongest code-supported explanation is that a fault in communication
with the mass flow controller (MFC) aborts a run, after which the software
intentionally attempts to place the system in its \code{off} state. This
sequence can look like a sudden shutdown even though the initiating event
was a serial communication or response error.

The standard reference method includes a monitored sample pulse lasting
31.25 minutes. During that pulse, the controller repeatedly reads the sample
flow over the serial connection. If communication stops responding, the
sequence can fail during the sampling step. The current logs contain examples
of \code{TimeoutError: Device timed out.} in the sampling step, including
failures around 3:13 AM and 4:11 AM. Other log entries show similar timeouts
at different times of day. This is evidence for a recurrent communication
failure path, but it does not establish what causes the communication failure.

\subsection{How an MFC response can become a run failure}

For a monitored pulse, the state handler calls \code{sample(display=False)}
while sample flow is active. The call chain is:
\begin{enumerate}
    \item The pulse monitor calls \code{sample(display=False)}.
    \item \code{sample()} calls \code{flowrate(position=0)} to query the
    sample MFC.
    \item \code{flowrate()} sends a request through the serial helper
    \code{uPy.echo()}.
    \item If the device does not complete its response prompt before the
    timeout, \code{uPy.echo()} raises \code{TimeoutError}. The recovery in
    \code{flowrate()} handles \code{SerialException}, but does not catch
    this timeout, so it propagates out of the sampling state.
\end{enumerate}

There is a second, distinct failure path if the serial helper receives an
empty response rather than timing out. The helper returns \code{None} for an
empty evaluated response. \code{flowrate()} appends that value to its result
list without checking it, and \code{sample(display=False)} can consequently
return \code{None}. The monitored pulse stores that result in its flow
measurement list. When the pulse completes, the code integrates the list to
calculate sampled volume; invalid or missing measurements can make this
calculation fail. This missing-value path is plausible from the source, but
the cited overnight log entries specifically report timeouts, not a
\code{None}-related integration error.

\subsection{Why a communication error can look like a shutdown}

When a method raises an exception, \code{run\_method()} catches it long
enough to write a failure record to \say{logs/method log.csv}. If the log
write succeeds, it then calls \code{state("off")} before raising an
exception to its caller. The \code{off} state commands valves and relays off,
then requests zero sample and backflush flow. A user watching the hardware
may therefore see a shutdown as the visible result of an earlier MFC, serial,
temperature, or GC error.

This shutdown attempt is not guaranteed to finish. It uses the same hardware
communication paths, including the MFC connection; if one of those calls
fails again, the off-state transition can itself raise an exception. Outputs
that were commanded before the failure may already be off, while later
commands may not have completed. Treat an apparent shutdown as a potentially
incomplete transition until the physical outputs have been checked.

This explanation applies if the Raspberry Pi process remained alive long
enough to handle an exception and command the attached system off. If the Pi
itself lost power, rebooted, or became unresponsive, this Python shutdown path
cannot by itself explain that event; investigate Pi power, undervoltage,
thermal conditions, storage, and operating-system logs as a separate failure
mode.

The sequence's default error policy is to re-raise a run failure, which ends
the sequence rather than automatically continuing. A message may be missed
if the program was launched without a visible terminal or its output was
redirected. Some code paths also redirect Python's standard output to the
null device while suppressing output; if an exception occurs before output
is restored, later diagnostic text may not reach the console. The method log
is therefore important, but it is not a substitute for preserving the
terminal output and checking the actual hardware state.

\subsection{Why it may happen after several runs}

The code does not schedule a shutdown for 2--4 AM, and the Python source
contains no time-of-day rule for stopping a run. An offline Raspberry Pi
does not need internet access for the serial communications or GPIO calls
shown here. A repeated long-running sequence does, however, perform many
hardware operations; intermittent communication faults may become more
likely to be encountered as the number of polls and runs increases. The
observed hour could also reflect when a run in a regularly spaced sequence
reaches its sampling or shutdown steps.

The project's \say{Day Run} sequence specifies a 45-minute minimum interval
between runs. The interval begins at the start of each run; it is not a
fixed time-of-day schedule. Whether that explains the observed timing depends
on how the sequence was started, whether continuous operation was enabled,
and the actual sequence file used at the time.

An archived method log also records \code{[Errno 28] No space left on
device} during a run. This confirms that storage exhaustion occurred at
least once historically, but it does not explain the later serial timeout
entries by itself. A full filesystem or exhausted inodes can prevent logging
and other writes, so available storage should still be checked when
investigating an apparently silent failure.

\subsection{Checks to perform after a failure}

Before restarting the sequence, verify the physical state of the pump,
valves, MFCs, traps, and GC; do not infer their state solely from the
program's last printed message. Then preserve the following evidence:
\begin{itemize}
    \item The full terminal output, including any traceback, from before and
    after the apparent shutdown.
    \item The corresponding final row in \say{logs/method log.csv}, including
    the method name, state name, timestamp, and error text.
    \item The Pi's available disk space and inodes (\code{df -h} and
    \code{df -i}), and whether the filesystem was mounted read-only.
    \item Kernel messages near the failure time (\code{dmesg}), looking for
    USB disconnects, USB resets, undervoltage, or power-related warnings.
    \item The selected sequence, its repetition and interval settings, and
    whether it was started with continuous operation or by a supervisor.
\end{itemize}

If it is safe to do so, record whether the MFC still responds to a direct
read after the failure and whether the serial device disappears from the
Pi. Also note whether the GC reports ready and whether temperature readings
remain available. These observations help distinguish an MFC-only
communication loss from a broader USB, power, Raspberry Pi, or instrument
shutdown.

\subsection{Information needed to confirm the hardware cause}

The source and logs support the software failure path above, but the
underlying physical cause remains undetermined. To narrow it down, record
the Raspberry Pi model and operating-system version; the power supply and
any USB hub used; which controllers share USB or power connections; and
whether the Pi itself stayed on when the pre-concentrator outputs shut down.
The MFC and valve-controller board types, their power arrangements, and
their serial connection details are also relevant.

\warn{If the controller appears to shut down during an unattended run,
confirm the physical state and restore a known safe condition before
restarting. A software exception or a log entry does not prove that every
output reached its intended safe state.}
